\section*{Chapter \ref{ch:g10:electron-shells} --- Electron Shells and the Periodic Table}

\begin{solution}{exo:g10:electron-shells:1}
C: $1\mathrm{s}^2\,2\mathrm{s}^2\,2\mathrm{p}^2$. N:
$1\mathrm{s}^2\,2\mathrm{s}^2\,2\mathrm{p}^3$. Mg:
$1\mathrm{s}^2\,2\mathrm{s}^2\,2\mathrm{p}^6\,3\mathrm{s}^2$.
\end{solution}

\begin{solution}{exo:g10:electron-shells:2}
Carbon 4, nitrogen 5, magnesium 2.
\end{solution}

\begin{solution}{exo:g10:electron-shells:3}
Outer \omterm{def:g10:electron-shells:configuration}{shell} $n = 3$: period 3. $2 + 3 = 5$ \omterm{def:g10:electron-shells:valence}{valence electrons}: \omterm{def:g9:periodic-table-first-look:period}{group} 15
(phosphorus).
\end{solution}

\begin{solution}{exo:g10:electron-shells:4}
\omterm{def:g10:electron-shells:family}{Alkali metals} (lithium, sodium, potassium); \omterm{def:g10:electron-shells:family}{halogens} (fluorine,
chlorine, bromine, iodine); \omterm{def:g10:electron-shells:family}{noble gases} (helium, neon, argon).
\end{solution}

\begin{solution}{exo:g10:electron-shells:5}
\omterm{def:g7:atoms-and-molecules:atom}{Atoms} tend to gain or lose \omterm{def:g9:inside-the-atom:nucleus}{electrons} so as to have eight \omterm{def:g9:inside-the-atom:nucleus}{electrons} in
their outer \omterm{def:g10:electron-shells:configuration}{shell}, like the nearest \omterm{def:g10:electron-shells:family}{noble gas}.
\end{solution}

\begin{solution}{exo:g10:electron-shells:6}
K, $2.8.8.1$, loses one: \ce{K+},
$1\mathrm{s}^2\,2\mathrm{s}^2\,2\mathrm{p}^6\,3\mathrm{s}^2\,3\mathrm{p}^6$
(argon). F, $2.7$, gains one: \ce{F-},
$1\mathrm{s}^2\,2\mathrm{s}^2\,2\mathrm{p}^6$ (neon).
\end{solution}

\begin{solution}{exo:g10:electron-shells:7}
$1\mathrm{s}^2\,2\mathrm{s}^2\,2\mathrm{p}^6\,3\mathrm{s}^2$: $Z = 12$,
magnesium.
\end{solution}

\begin{solution}{exo:g10:electron-shells:8}
\ce{Mg^{2+}} and \ce{Cl-}: \ce{MgCl2}. \ce{Al^{3+}} and \ce{O^{2-}}:
\ce{Al2O3}.
\end{solution}

\begin{solution}{exo:g10:electron-shells:9}
Sulfur (6 \omterm{def:g10:electron-shells:valence}{valence electrons}, $2.8.6$), in the same group, 16.
\end{solution}

\begin{solution}{exo:g10:electron-shells:10}
Its outer \omterm{def:g10:electron-shells:configuration}{shell} is already full (eight \omterm{def:g9:inside-the-atom:nucleus}{electrons}): it has nothing to gain
by losing or gaining \omterm{def:g9:inside-the-atom:nucleus}{electrons}.
\end{solution}

\begin{solution}{exo:g10:electron-shells:11}
Argon has 18 \omterm{def:g9:inside-the-atom:nucleus}{electrons}; \ce{X^{2+}} has lost 2, so X has 20: calcium.
\end{solution}

\begin{solution}{exo:g10:electron-shells:12}
Helium, $1\mathrm{s}^2$, has a full outer \omterm{def:g10:electron-shells:configuration}{shell} (\omterm{def:g10:electron-shells:configuration}{shell} 1 holds only two
\omterm{def:g9:inside-the-atom:nucleus}{electrons}), like the other \omterm{def:g10:electron-shells:family}{noble gases}; it does not behave like the \omterm{def:g9:periodic-table-first-look:period}{group} 2
\omterm{def:g9:periodic-table-first-look:metal}{metals}.
\end{solution}

\begin{solution}{exo:g10:electron-shells:13}
Its \omterm{def:g10:electron-shells:valence}{valence electron} is in \omterm{def:g10:electron-shells:configuration}{shell} 4, further from the \omterm{def:g9:inside-the-atom:nucleus}{nucleus} than sodium's
in \omterm{def:g10:electron-shells:configuration}{shell} 3: less tightly held, it is lost more easily.
\end{solution}

\begin{solution}{exo:g10:electron-shells:14}
All four have 18 \omterm{def:g9:inside-the-atom:nucleus}{electrons} and the configuration of argon:
\[ 1\mathrm{s}^2\,2\mathrm{s}^2\,2\mathrm{p}^6\,3\mathrm{s}^2\,3\mathrm{p}^6. \]
\end{solution}

\begin{solution}{exo:g10:electron-shells:15}
It gains 3 \omterm{def:g9:inside-the-atom:nucleus}{electrons} to reach an octet: it has 5 \omterm{def:g10:electron-shells:valence}{valence electrons}, \omterm{def:g9:periodic-table-first-look:period}{group}
15. In period 3, it is phosphorus:
\[ 1\mathrm{s}^2\,2\mathrm{s}^2\,2\mathrm{p}^6\,3\mathrm{s}^2\,3\mathrm{p}^3. \]
\end{solution}

\begin{solution}{pb:g10:electron-shells:1}
\textbf{1.} Aluminium: 13 \omterm{def:g9:inside-the-atom:proton}{protons}, 13 \omterm{def:g9:inside-the-atom:nucleus}{electrons}. Oxygen: 8 \omterm{def:g9:inside-the-atom:proton}{protons}, 8
\omterm{def:g9:inside-the-atom:nucleus}{electrons}.

\textbf{2.} Al: $1\mathrm{s}^2\,2\mathrm{s}^2\,2\mathrm{p}^6\,3\mathrm{s}^2\,3\mathrm{p}^1$.
O: $1\mathrm{s}^2\,2\mathrm{s}^2\,2\mathrm{p}^4$.

\textbf{3.} Al: 3 \omterm{def:g10:electron-shells:valence}{valence electrons}, period 3, \omterm{def:g9:periodic-table-first-look:period}{group} 13. O: 6, period 2,
\omterm{def:g9:periodic-table-first-look:period}{group} 16.

\textbf{4.} Aluminium is a \omterm{def:g9:periodic-table-first-look:metal}{metal}, oxygen a \omterm{def:g9:periodic-table-first-look:metal}{non-metal}.

\textbf{5.} \ce{Al^{3+}}: $1\mathrm{s}^2\,2\mathrm{s}^2\,2\mathrm{p}^6$.

\textbf{6.} \ce{O^{2-}}: $1\mathrm{s}^2\,2\mathrm{s}^2\,2\mathrm{p}^6$.

\textbf{7.} Neon.

\textbf{8.} Aluminium loses 3, oxygen gains 2.

\textbf{9.} $2 \times (+3) + 3 \times (-2) = 0$: \ce{Al2O3}.

\textbf{10.} Lost: $2 \times 3 = 6$. Gained: $3 \times 2 = 6$.

\textbf{11.} \omterm{def:g9:inside-the-atom:nucleus}{Electrons} are not created or destroyed: those lost by
aluminium are exactly those gained by oxygen, and the compound is
neutral.

\textbf{12.} It carries the same charge, $+3$: the charges of the crystal
still balance.

\textbf{13.} $2 \times 27 + 3 \times 16 = 102$ \omterm{def:g9:inside-the-atom:proton}{nucleons}.

\textbf{14.} $102 \times 1.67 \times 10^{-27} = \qty{1.70e-25}{kg}$.

\textbf{15.} $10^{-3} / 1.70 \times 10^{-25} \approx 5.87 \times 10^{21}$
formula units.

\textbf{16.} $2 \times 5.87 \times 10^{21} = 1.17 \times 10^{22}$
\ce{Al^{3+}} \omterm{def:g9:ions:ion}{ions}; $3 \times 5.87 \times 10^{21} = 1.76 \times 10^{22}$
\ce{O^{2-}} \omterm{def:g9:ions:ion}{ions}.

\textbf{17.} An \omterm{def:g9:inside-the-atom:nucleus}{electron} is about 1836 times lighter than a \omterm{def:g9:inside-the-atom:proton}{nucleon}; the
50 \omterm{def:g9:inside-the-atom:nucleus}{electrons} of a formula unit weigh less than three hundredths of one
\omterm{def:g9:inside-the-atom:proton}{nucleon}.

\textbf{18.} $6 \times 5.87 \times 10^{21} \approx 3.5 \times 10^{22}$
\omterm{def:g9:inside-the-atom:nucleus}{electrons} transferred for \qty{1}{g} of corundum.
\end{solution}
